\documentclass[12pt]{article}

% --- Packages ---
\usepackage[utf8]{inputenc}
\usepackage[margin=1in]{geometry}
\usepackage{titlesec}
\usepackage{tabularx}
\usepackage{hyperref}
\usepackage{enumitem}
\usepackage{booktabs}

% --- Setup ---
\hypersetup{colorlinks=true, linkcolor=blue, urlcolor=cyan}

\begin{document}

% --- i. Cover Page ---
\begin{titlepage}
    \centering
    \vspace*{2cm}
    {\Huge \textbf{Software Requirement Specification}} \\
    \vspace{0.5cm}
    {\Large \textbf{JPL Lunar Detection Pipeline}} \\
    \vfill
    \textbf{Version:} 1.1 \\
    \textbf{Prepared by:} \\
     Ricky Chao, George Malanche, Kevin Mendoza, Connor Moy \\
    \vspace{1cm}
    \textbf{Sponsored by:} NASA JPL \\
    \textbf{Date:} May 14, 2026
\end{titlepage}

% --- ii. Table of Contents ---
\newpage
\tableofcontents
\newpage

% --- iii. Versions Table ---
\section*{Revision History}
\begin{tabularx}{\textwidth}{|l|l|l|X|}
\hline
\textbf{Revision} & \textbf{Date} & \textbf{Author} & \textbf{Description} \\ \hline
1.0 & 2026-05-04 & Project Team & Initial SRS Draft \\ \hline
1.1 & 2026-05-13 & Project Team & Update Sections 1.1 - 4 \\ \hline
\end{tabularx}

% --- iv. Sections ---

% 1. Introduction
\section{Introduction}
\subsection{Purpose}
The purpose of this SRS is to define the functional and non-functional requirements for the JPL Lunar Detection Pipeline. This uses an automated system for the identification and labeling of lunar surface features (craters, rocks, pits) to assist NASA scientists.

\subsection{Intended Audience}
This document is intended for NASA JPL stakeholders, project developers, and the verification/validation team.

\subsection{Overview}
The pipeline ingests raw LROC imagery, processes it through machine learning models (YOLO and Detectron2), and outputs geolocated data compatible with planetary mapping standards.

% 2. External Interface Requirements
\section{External Interface Requirements}
\subsection{User Interface}
The system will provide a Command Line Interface (CLI) for batch processing and a set of API endpoints for integration into existing JPL workflows. There is no requirement for a high-fidelity Graphical User Interface in this version.

\subsection{Software Interfaces}
\begin{itemize}
    \item \textbf{NASA PDS4:} Interface for retrieving raw planetary data and XML labels.
    \item \textbf{SPICE Toolkit:} Used for transforming image pixel coordinates into lunar geographic coordinates (Latitude/Longitude).
    \item \textbf{AWS S3:} Future interface for cloud-based storage of training datasets and model weights.
\end{itemize}

% 3. Legal and Ethical Considerations
\section{Legal and Ethical Considerations}
\subsection{User Data \& Storing}
The system primarily processes public-domain planetary data from NASA. However, any user-contributed labels or metadata stored in the pipeline must be handled according to NASA security protocols to ensure data integrity and prevent unauthorized modification.

\subsection{Legal Compliance}
All software components must adhere to open-source licensing (e.g., Apache 2.0 or MIT) to ensure the pipeline can be maintained and shared within the scientific community without proprietary restrictions.

\subsection{Ethical Dilemmas}
\begin{itemize}
    \item \textbf{Scientific Accuracy:} There is an ethical responsibility to ensure high precision in detections. False positives in crater labeling could lead to incorrect scientific conclusions regarding lunar age or landing site safety.
    \item \textbf{Data Provenance:} All processed images must maintain metadata linking them back to the original LROC source to ensure reproducibility in scientific research.
\end{itemize}

% --- v. Glossary ---
\newpage
\section{Glossary}
\begin{tabularx}{\textwidth}{l|X}
\toprule
\textbf{Acronym} & \textbf{Definition} \\ \midrule
LROC & Lunar Reconnaissance Orbiter Camera. \\
PDS4 & Planetary Data System 4; NASA's standard for planetary data labels. \\
SPICE & Spacecraft Planet Instrument C-matrix Events; NASA navigation toolkit. \\
YOLO & You Only Look Once; a real-time object detection model. \\ \bottomrule
\end{tabularx}

\end{document}